\documentclass[activite]{thedocument}
\startdatakeys
\source{}
\cycle{}
\competencesDuSocle{}
\connaissancesRequises{}
\competencesTravaillees{}
\enddatakeys
\begin{document}
\begin{enumerate}
    \item Selon vous qu'est-ce que les fractions suivantes ont en commun~?\vskip20pt
    \begin{center}
        \begin{tikzpicture}[scale=0.8]
            \coordinate(A1)at(-2,1);
            \coordinate(A2)at(0,0);
            \coordinate(A3)at(2.5,-2);
            \coordinate(A4)at(3.5,1);
            \coordinate(A5)at(4,-0.75);
            \draw node at(A1){$\dfrac{58}{116}$};
            \draw node at(A2){$\dfrac{2}{4}$};
            \draw node at(A3){$\dfrac{9}{18}$};
            \draw node at(A4){$\dfrac{5}{10}$};
            \draw node at(A5){$\dfrac{23}{46}$};
            \draw[dotted](-3,-3)rectangle(5,2);
        \end{tikzpicture}
    \end{center}\dotedline\dotedline\dotedline
    \item Combien de fractions similaires à celles-ci pouvez-vous écrire~?
    \dotedline\dotedline
    \item Quelle est la plus \guillemets{simple} de toutes les fractions
    de cette catégorie~? Pourquoi~?
    \dotedline\dotedline\dotedline
    \item Inventer une autre catégorie.
    \dotedline\dotedline\dotedline
\end{enumerate}


\end{document}